% Ad Familiares 2.13 (ad-familiares-02-13)
% Source: https://raw.githubusercontent.com/PerseusDL/canonical-latinLit/master/data/phi0474/phi056/phi0474.phi056.perseus-lat1.xml
% Fetched by scripts/fetch_latin.py

% Dateline (Perseus): Scr. Laudiceae inter K. et Non. Maias a. 704 (50) .

\ciceroLetterOpener{M. CICERO IMP. S. D. M. CAELIO AEDILI CVRVLI.}

\ciceroSection{1} raras tuas quidem (fortasse enim non perferuntur), sed suavis accipio litteras; vel quas proxime acceperam, quam prudentis, quam multi et offici et consili! etsi omnia sic constitueram mihi agenda, ut tu admonebas tamen conio firmantur nostra consilia, cum sentimus prudentibus fideliterque suadentibus idem videri.

\ciceroSection{2} ego Appium, ut saepe tecum locutus sum, valde diligo meque ab eo diligi statim coeptum esse, ut simultatem deposuimus, sensi; nam et honorificus in me cos. fuit et suavis amicus et studiosus studiorum etiam meorum; mea vero officia ei non defuisse tu es testis, quoi iam kwmiko\s martu/s ut opinor, accedit Phania, et me hercule etiam pluris eum feci, quod te amari ab eo sensi. iam me Pompei totum esse scis; Brutum a me amari intellegis. quid est causae cur mihi non in optatis sit complecti hominem florentem aetate, opibus, honoribus, ingenio, liberis, propinquis, adfinibus, amicis, conlegam meum praesertim et in ipsa collegi laude et scientia studiosum mei? haec eo pluribus scripsi, quod non nihil significabant tuae litterae subdubitare, qua essem erga illum voluntate. credo te audisse aliquid. falsum est, mihi crede, si quid audisti. genus institutorum et rationum mearum dissimilitudinem non nullam habet cum illius administratione provinciae; ex eo quidam suspicati fortasse sunt animorum contentione, non opinionum dissensione me ab eo discrepare; nihil autem feci umquam neque dixi, quod contra illius existimationem esse vellem, post hoc negotium autem et temeritatem nostri Dolabellae deprecatorem me pro illius periculo praebeo.

\ciceroSection{3} erat in eadem epistula 'veternus civitatis.' gaudebam sane et congelasse nostrum amicum laetabar otio. extrema pagella pupugit me tuo chirographo. quid ais? Caesarem nunc defendit 'Curio? quis hoc putaret praeter me? nam, ita vivam, putavi. di immortales, quam ego risum nostrum desidero! mihi erat in animo, quoniam iuris dictionem confeceram, civitates locupletaram, publicanis etiam superioris lustri reliqua sine sociorum ulla querela conservaram, privatis, summis infimis, fueram iucundus, proficisci in Ciliciam Nonis Mais et, cum prima aestiva attigissem militemque conlocassem, decedere ex senatus consulto. cupio te aedilem videre, miroque desiderio me urbs adficit et omnes mei tuque in primis.
